% ============================================================
\chapter{$L^p$ Spaces}
\label{ch:lp_spaces}
% ============================================================

The $L^p$ spaces are the most important Banach spaces in analysis. Born from the Lebesgue integral, they provide the natural framework for Fourier analysis, partial differential equations, probability theory, and machine learning. Frigyes Riesz, in the 1910s, was the first to study these spaces systematically, establishing with Ernst Fischer the celebrated completeness theorem for $L^2$. The H\"older and Minkowski inequalities structure the theory; the duality $(L^p)^* = L^q$ (for $1/p + 1/q = 1$) crowns it. As for $L^2$, endowed with its inner product, it is the only one that is a Hilbert space, and it is upon $L^2$ that all of spectral theory rests.

\section{Definitions}

\begin{definition}[$L^p$ space]
Let $(X, \mathcal{A}, \mu)$ be a measure space and $1 \leq p < \infty$. The space
$\mathcal{L}^p(\mu)$ is the set of measurable functions $f : X \to \R$ (or $\C$)
such that:
\[
\int_X \abs{f}^p\, d\mu < \infty.
\]
The \emph{$L^p$ norm} (or \emph{semi-norm} on $\mathcal{L}^p$) is:
\[
\norm{f}_p = \left(\int_X \abs{f}^p\, d\mu\right)^{1/p}.
\]
The space $L^p(\mu) = \mathcal{L}^p(\mu) / \sim$ is the quotient by the equivalence
relation $f \sim g \iff f = g$ $\mu$-a.e.
\end{definition}

\begin{definition}[$L^\infty$ space]
The space $L^\infty(\mu)$ is the set of equivalence classes of measurable
essentially bounded functions, equipped with the norm:
\[
\norm{f}_\infty = \inf\{M \geq 0 : \abs{f} \leq M \text{ $\mu$-a.e.}\}
= \esssup \abs{f}.
\]
\end{definition}

\begin{remark}
One works with equivalence classes (modulo a.e. equality) so that $\norm{\cdot}_p$
is a genuine norm (not merely a semi-norm).
\end{remark}

\section{H\"older's inequality}

\begin{definition}[Conjugate exponents]
Two reals $p, q \in [1, +\infty]$ are \emph{conjugate} (or \emph{dual}) if:
\[
\frac{1}{p} + \frac{1}{q} = 1.
\]
By convention, $q = \infty$ if $p = 1$ and $q = 1$ if $p = \infty$.
\end{definition}

\begin{lemma}[Young's inequality]
For $a, b \geq 0$ and conjugate $p, q > 1$:
\[
ab \leq \frac{a^p}{p} + \frac{b^q}{q}.
\]
\end{lemma}

\begin{proof}
The function $t \mapsto \log t$ is concave, so for $\alpha = 1/p$ and
$\beta = 1/q$:
\[
\log(\alpha a^p + \beta b^q) \geq \alpha \log(a^p) + \beta \log(b^q)
= \log(a) + \log(b) = \log(ab).
\]
Exponentiating: $ab \leq \frac{a^p}{p} + \frac{b^q}{q}$.
\end{proof}

\begin{theorem}[H\"older's inequality]
\label{thm:holder}
Let $p, q \in [1, +\infty]$ be conjugate. If $f \in L^p(\mu)$ and
$g \in L^q(\mu)$, then $fg \in L^1(\mu)$ and:
\[
\norm{fg}_1 = \int_X \abs{fg}\, d\mu \leq \norm{f}_p \cdot \norm{g}_q.
\]
\end{theorem}

\begin{proof}
The case $p = 1$, $q = \infty$ is immediate: $\abs{fg} \leq \abs{f} \cdot
\norm{g}_\infty$ a.e.

For $1 < p < \infty$, we may assume $\norm{f}_p = \norm{g}_q = 1$ (otherwise
normalize). By Young's inequality applied to $a = \abs{f(x)}$ and
$b = \abs{g(x)}$:
\[
\abs{f(x) g(x)} \leq \frac{\abs{f(x)}^p}{p} + \frac{\abs{g(x)}^q}{q}.
\]
Integrating:
\[
\int \abs{fg}\, d\mu \leq \frac{\norm{f}_p^p}{p} + \frac{\norm{g}_q^q}{q}
= \frac{1}{p} + \frac{1}{q} = 1 = \norm{f}_p \cdot \norm{g}_q. \qedhere
\]
\end{proof}

\begin{corollary}[Cauchy--Schwarz inequality]
For $p = q = 2$: $\int \abs{fg}\, d\mu \leq \norm{f}_2 \cdot \norm{g}_2$.
\end{corollary}

\section{Minkowski's inequality}

\begin{theorem}[Minkowski's inequality]
\label{thm:minkowski}
Let $1 \leq p \leq \infty$. For $f, g \in L^p(\mu)$:
\[
\norm{f + g}_p \leq \norm{f}_p + \norm{g}_p.
\]
That is, $\norm{\cdot}_p$ is a norm on $L^p(\mu)$.
\end{theorem}

\begin{proof}
The case $p = 1$ is the triangle inequality. The case $p = \infty$ is clear.

For $1 < p < \infty$:
\begin{align*}
\int \abs{f + g}^p\, d\mu &\leq \int \abs{f + g}^{p-1}(\abs{f} + \abs{g})\, d\mu
\\
&= \int \abs{f + g}^{p-1}\abs{f}\, d\mu + \int \abs{f + g}^{p-1}\abs{g}\, d\mu.
\end{align*}
Apply H\"older to each term with $q = p/(p-1)$:
\[
\int \abs{f + g}^{p-1}\abs{f}\, d\mu \leq
\left(\int \abs{f + g}^p\, d\mu\right)^{1/q} \norm{f}_p.
\]
Thus $\norm{f + g}_p^p \leq \norm{f + g}_p^{p/q}(\norm{f}_p + \norm{g}_p)$.
Dividing by $\norm{f + g}_p^{p/q}$ (if nonzero) and noting $p - p/q = 1$:
$\norm{f + g}_p \leq \norm{f}_p + \norm{g}_p$.
\end{proof}

\begin{keyformulas}[title=\textbf{Fundamental inequalities}]
\[
\text{H\"older: } \norm{fg}_1 \leq \norm{f}_p \norm{g}_q, \qquad
\text{Minkowski: } \norm{f + g}_p \leq \norm{f}_p + \norm{g}_p.
\]
\end{keyformulas}

\section{Completeness: $L^p$ is a Banach space}

\begin{theorem}[Riesz--Fischer]
\label{thm:riesz_fischer}
For every $1 \leq p \leq \infty$, the space $(L^p(\mu), \norm{\cdot}_p)$ is a
Banach space (complete normed vector space).
\end{theorem}

\begin{proof}
Let $(f_n)$ be a Cauchy sequence in $L^p$. It suffices to show it has a convergent
subsequence (every Cauchy sequence with a convergent subsequence converges).

Choose $n_k$ so that $\norm{f_{n_{k+1}} - f_{n_k}}_p < 2^{-k}$. Set:
\[
g_N = \sum_{k=1}^{N} \abs{f_{n_{k+1}} - f_{n_k}}, \qquad
g = \sum_{k=1}^{\infty} \abs{f_{n_{k+1}} - f_{n_k}}.
\]
By Minkowski, $\norm{g_N}_p \leq \sum_{k=1}^N 2^{-k} \leq 1$. By monotone
convergence, $\norm{g}_p \leq 1$, so $g < \infty$ a.e. This means the telescoping
series $f_{n_1} + \sum_{k=1}^{\infty}(f_{n_{k+1}} - f_{n_k})$ converges absolutely
a.e. Let $f$ be its sum (a.e.). Then $f_{n_k} \to f$ a.e. and
$\abs{f_{n_k} - f} \leq 2g \in L^p$, so by dominated convergence,
$\norm{f_{n_k} - f}_p \to 0$.
\end{proof}

\section{Inclusion relations between $L^p$ spaces}

\begin{proposition}
If $\mu(X) < \infty$ and $1 \leq p \leq q \leq \infty$, then
$L^q(\mu) \subset L^p(\mu)$ and:
\[
\norm{f}_p \leq \mu(X)^{1/p - 1/q} \norm{f}_q.
\]
\end{proposition}

\begin{proof}
Apply H\"older with $r = q/p \geq 1$ to $\abs{f}^p$ and $1$:
$\int \abs{f}^p\, d\mu \leq \left(\int \abs{f}^q\, d\mu\right)^{p/q}
\mu(X)^{1 - p/q}$.
\end{proof}

\begin{warning}[title=\textbf{No inclusion in general}]
If $\mu$ is not finite, there is generally no inclusion between the $L^p$ spaces.
For example, on $(\R, \lambda)$: $f(x) = x^{-1/2}\mathbf{1}_{(0,1)} \in L^1
\setminus L^2$ and $g(x) = (1+x)^{-1} \in L^2(\R) \setminus L^1(\R)$.
\end{warning}

\section{Density and separability}

\begin{theorem}
For $1 \leq p < \infty$:
\begin{enumerate}
    \item Integrable simple functions are dense in $L^p(\mu)$.
    \item If $X = \R^n$ and $\mu = \lambda^n$, the continuous functions with compact
    support $C_c(\R^n)$ are dense in $L^p(\R^n)$.
    \item $L^p(\R^n)$ is separable for $1 \leq p < \infty$.
\end{enumerate}
\end{theorem}

\begin{remark}
$L^\infty$ is generally \emph{not} separable.
\end{remark}

\section{Duality of $L^p$ spaces}

\begin{theorem}[$L^p$--$L^q$ duality]
\label{thm:lp_duality}
Let $1 \leq p < \infty$ and $q$ the conjugate exponent. For every continuous linear
functional $\ell : L^p(\mu) \to \R$, there exists a unique $g \in L^q(\mu)$ such
that:
\[
\ell(f) = \int_X fg\, d\mu, \quad \forall f \in L^p(\mu).
\]
Moreover, $\norm{\ell} = \norm{g}_q$. In other words,
$(L^p(\mu))^* \cong L^q(\mu)$ (isometrically).
\end{theorem}

\begin{remark}
This theorem requires $\mu$ to be $\sigma$-finite. For $p = 1$,
$(L^1)^* \cong L^\infty$. For $p = \infty$, $(L^\infty)^* \supsetneq L^1$ in
general.
\end{remark}

\section{$L^2$ as a Hilbert space}

\begin{definition}[$L^2$ inner product]
The space $L^2(\mu)$ is equipped with the inner product:
\[
\inner{f}{g} = \int_X f \overline{g}\, d\mu.
\]
It is complete for the associated norm $\norm{f}_2 = \sqrt{\inner{f}{f}}$, hence
it is a Hilbert space.
\end{definition}

\begin{theorem}[Orthogonal projection]
If $V$ is a closed subspace of $L^2(\mu)$ and $f \in L^2(\mu)$, there exists a
unique $\hat{f} \in V$ such that:
\[
\norm{f - \hat{f}}_2 = \inf_{g \in V} \norm{f - g}_2.
\]
Moreover, $f - \hat{f} \perp V$ (i.e.\ $\inner{f - \hat{f}}{g} = 0$ for all
$g \in V$).
\end{theorem}

\section{Complementary inequalities}

\begin{proposition}[Jensen's inequality]
\label{prop:jensen}
Let $(X, \mathcal{A}, \mu)$ be a probability space ($\mu(X) = 1$),
$f \in L^1(\mu)$, and $\varphi : \R \to \R$ convex. If $\varphi \circ f$ is
integrable, then:
\[
\varphi\!\left(\int_X f\, d\mu\right) \leq \int_X \varphi \circ f\, d\mu.
\]
\end{proposition}

\begin{proposition}[Interpolation inequality]
If $f \in L^{p_1} \cap L^{p_2}$ with $1 \leq p_1 < p_2 \leq \infty$, then
$f \in L^p$ for every $p \in [p_1, p_2]$ and:
\[
\norm{f}_p \leq \norm{f}_{p_1}^{\theta} \norm{f}_{p_2}^{1-\theta}, \quad
\text{where } \frac{1}{p} = \frac{\theta}{p_1} + \frac{1-\theta}{p_2}.
\]
\end{proposition}

\section{Exercises}

\begin{exercise}
Show that for $1 \leq p < q < \infty$, $\ell^p \subset \ell^q$ (where
$\ell^p = L^p(\N, \mathcal{P}(\N), \#)$ with $\#$ the counting measure). This is
the reverse of the inclusion for finite measure spaces.
\end{exercise}

\begin{exercise}
Show that $L^2([0,1])$ is a separable Hilbert space and that
$(e_n)_{n \in \Z}$ defined by $e_n(x) = e^{2\pi inx}$ is an orthonormal basis.
\end{exercise}

\begin{exercise}
Prove the integral form of Minkowski's inequality: if $f \geq 0$ is measurable on
$X \times Y$ and $1 \leq p < \infty$, then
\[
\left(\int_X \left(\int_Y f(x,y)\, d\nu(y)\right)^p d\mu(x)\right)^{1/p}
\leq \int_Y \left(\int_X f(x,y)^p\, d\mu(x)\right)^{1/p} d\nu(y).
\]
\end{exercise}

\begin{exercise}
Let $(f_n) \subset L^p(\mu)$ with $\sum_n \norm{f_n}_p < \infty$. Show that
$\sum_n f_n$ converges a.e. and in $L^p$.
\end{exercise}

\begin{exercise}
Show that if $f_n \to f$ in $L^p$ and $f_n \to g$ a.e., then $f = g$ a.e.
\end{exercise}
